\documentclass{article}
\usepackage[en]{homework}

\config{Algebra-0}{2026 Spring}{6}

\begin{document}


\begin{problem}[5A-29]
Suppose $T \in \mathcal{L}(\R^3)$ and $-4$, $5$, and $\sqrt{7}$ are eigenvalues of $T$. Prove that there exists $x \in \R^3$ such that
\[
Tx - 9x = (-4, 5, \sqrt{7}).
\]
\end{problem}

\begin{proof}
  Since $-4$, $5$, and $\sqrt{7}$ are eigenvalues of $T\in \mathcal{L}(\R^3)$, we can infer that \( 9 \) is not an eigenvalue, hence \( \ker (T-9I)=\left\{ 0 \right\} \). Therefore, \( T-9I \) is injective. Since \( \R^3 \) is finite-dimensional, injectivity implies surjectivity. Thus, for any vector \( v \) in \( \R^3 \), there exists a vector \( x \in \R^3 \) such that \( Tx-9x=v \). Take \( v=(-4,5,\sqrt{7}) \), then the problem has been proved.
\end{proof}



\begin{problem}[5A-30]
Suppose $T \in \mathcal{L}(V)$ and
\[
(T - 2I)(T - 3I)(T - 4I) = 0.
\]
Suppose $\lambda$ is an eigenvalue of $T$. Prove that $\lambda = 2$ or $\lambda = 3$ or $\lambda = 4$.
\end{problem}

\begin{proof}
  Take \( v\neq0 \) such that \( Tv=\lambda v \). Then
  \[ 0=(T-2I)(T-3I)(T-4I)v=(\lambda-2)(\lambda-3)(\lambda-4)v\implies\lambda\in \left\{ 2, 3, 4 \right\}. \]
\end{proof}



\begin{problem}[5A-38]
Suppose $V$ is finite-dimensional, $T \in \mathcal{L}(V)$, and $U$ is a subspace of $V$ invariant under $T$. The quotient operator $T/U \in \mathcal{L}(V/U)$ is defined by
\[
(T/U)(v + U) = Tv + U
\]
for each $v \in V$.
\begin{enumerate}[label=(\alph*)]
	\item Show that the definition of $T/U$ makes sense (which requires using the condition that $U$ is invariant under $T$) and show that $T/U$ is an operator on $V/U$.
	\item Show that each eigenvalue of $T/U$ is an eigenvalue of $T$.
\end{enumerate}
\end{problem}

\begin{proof}
  (a) For any \( v_1,v_2\in V \) such that \( v_1+U=v_2+U \), we know that \( u:=v_1-v_2\in U \). Then
  \begin{align*}
    (T/U)(v_1+U) &=T(v_1)+U = T(v_2+u)+U = T(v_2)+Tu+U\\
    &= T(v_2)+U = (T/U)(v_2+U).
  \end{align*}
  Then \( T/U \) is well-defined.\par
  (b) Take an eigenvalue and corresponding eigenvector of \( T/U \), say \( \lambda \) and \( v+U \). Note that \( v+U \neq U \), so \( v \notin U \). Then \( Tv+U=(T/U)(v+U)=\lambda(v+U)=\lambda v+U \), which implies that \( Tv=\lambda v+u \) for some \( u\in U \). \par
  If \( \lambda \) is not an eigenvector of \( T \), then \( T-\lambda I \) is injective in \( V \), as well as in \( U \) (because \( U \) is T-invariant). Then there exists \( u_0\in U \) such that \( Tu_0-\lambda u_0=u\in U \). \par
  Set \( v'=v-u_0 \). Since \( v\notin U \), we have \( v'\neq 0 \). And by the definition of \( u_0 \), we have
  \[ Tv'=Tv-Tu_0=\lambda v+u-(\lambda u_0+u)=\lambda (v-u_0)=\lambda v'. \]
  This contradicts the assumption that \( \lambda \) is not an eigenvalue of \( T \). Therefore, \( \lambda \) is an eigenvalue of \( T \).
\end{proof}



\begin{problem}[5A-42]
Define $T \in \mathcal{L}(\F^n)$ by
\[
T(x_1, x_2, x_3, \ldots, x_n) = (x_1, 2x_2, 3x_3, \ldots, nx_n).
\]
\begin{enumerate}[label=(\alph*)]
	\item Find all eigenvalues and eigenvectors of $T$.
	\item Find all subspaces of $\F^n$ that are invariant under $T$.
\end{enumerate}
\end{problem}

\begin{solution}
  (a) The representation of \( T \) under the standard basis is \( \operatorname{diag}(1,2,\dots,n) \), then all eigenvalues and eigenvectors of \( T \) are
  \[ k\text{ and }(0,\dots,0,\overset{(k\text{ th})}{1},0,\dots,0), \quad k=1,2,\dots,n. \]\par
  (b) We will prove that the invariant subspaces of \( T \) are exactly those spanned by a subset of the standard basis, i.e. those of the form
  \[ V_S:=\left\{ a_1,\dots,a_n \in\F^n:\, a_i=0 \text{ for all } i\notin S \right\},\text{ where } S \subseteq \{1,2,\dots,n\}. \]\par
  Obviously, \( V_S \) is invariant under \( T \) for any \( S \subseteq \{1,2,\dots,n\} \). \par
  Conversely, take an invariant subspace \( U \) of \( T \). For any \( v=(a_1,a_2,\dots,a_n) \), Define 
  \[ N(v)=\left\{ k\in\left\{ 1,2,\dots,n \right\}:a_k\neq 0 \right\}, \]
  \[ v_-=\min\left\{ \abs{a_i}:i\in N(v) \right\},\quad v_+=\max\left\{ \abs{a_i}:i\in N(v) \right\}. \]
  Then, let \( S=\bigcup_{u\in U}^{}N(u) \). We will show that \( U=V_S \). Obviously \( U\subset V_S \), so it suffices to prove that \( V_S\subset U \).\par
  First, we will show that there exists \( x\in U \) such that \( N(x)=S \). Denote \( S \) by \( S:=\left\{ s_1<\dots<s_m \right\} \). For each \( s_i \), there exists \( u_i\in U \) such that \( s_i\in N(u_i) \). We will construct \( x \) inductively by following steps:
  \begin{itemize}
    \item Take \( x_1=u_1 \). Then the \( s_1 \)-th component of \( x_1 \) is nonzero.
    \item Once we have chosen \( x_k \) such that \( s_i \)-th component of \( x_k \) is nonzero for all \( 1\le i\le k \). If the \( s_{k+1} \)-th component of \( x_k \) is nonzero, then take \( x_{k+1}=x_k \); otherwise let
    \[ x_{k+1} = x_k + \frac{(x_{k})_-}{2(u_{k+1})_+} u_{k+1}. \]
    The \( s_{k+1} \)-th component of \( x_{k+1} \) is nonzero due to the choice of \( u_{k+1} \); for any \( s_{i(\le k)} \)-th component of \( x_{k+1} \), we have
    \[ \abs{x_{k+1}^{(s_i)}}
    =\abs{x_{k}^{(s_i)}+\frac{(x_{k})_-}{2(u_{k+1})_+}u_{k+1}^{s_i}}
    \ge \abs{x_{k}^{(s_i)}}-\frac{(x_{k})_-}{2(u_{k+1})_+}\abs{u_{k+1}^{(s_i)}}
    \ge \abs{x_{k}^{(s_i)}}-\frac{1}{2}(x_k)_->0. \]
    Then the \( s_i \)-th component of \( x_{k+1} \) is nonzero. \par
    Hence, we have constructed \( x_{k+1} \) such that the \( s_i \)-th (\( 1\le i\le k+1 \)) components are nonzero.
    \item Finally, we have \( x=x_m \) such that \( N(x)=S \).
  \end{itemize}\par
  Next, we will show that \( V_S\subset\sspan\left\{ x,Tx,\dots,T^{m-1}x \right\}\subset U \).\par
  Denote \( a_i:=x^{(s_i)} \) for any \( 1\le i\le m \). Then, for any \( v\in V \) (\( v^{(s_i)}=b_i \)), the equation
  \[\left\{
    \begin{matrix}
      \lambda_1a_1&+\lambda_2a_1&+\lambda_3a_1&+\dots&+\lambda_m a_1&=b_1,\\
      \lambda_1a_2&+\lambda_2 2a_2&+\lambda_3 2^2a_2&+\dots&+\lambda_m 2^{m-1}a_2&=b_2,\\
      \lambda_1a_3&+\lambda_2 3a_3&+\lambda_3 3^2a_3&+\dots&+\lambda_m 3^{m-1}a_3&=b_3,\\
      \dots\\
      \lambda_1a_m&+\lambda_2 ma_m&+\lambda_m m^2 a_m&+\dots&+\lambda_m m^{m-1}a_m&=b_m
    \end{matrix}
  \right.\]
  has a solution \( (\lambda_1,\dots,\lambda_m) \) because the coefficient matrix is a Vandermonde matrix. Hence, 
  \[ v=\sum_{i=1}^{m} \lambda_i T^{i-1}x \in \sspan\left\{ x,Tx,\dots,T^{m-1}x \right\}\subset U. \]
  Therefore, \( V_S\subset U \). \par
  In conclusion, the invariant subspaces of \( T \) are exactly those of the form \( V_S \) for some \( S\subseteq \{1,2,\dots,n\} \).
\end{solution}



\begin{problem}[5A-43]
Suppose that $V$ is finite-dimensional, $\dim V > 1$, and $T \in \mathcal{L}(V)$. Prove that
\[
\{p(T) : p \in \mathcal{P}(\F)\} \neq \mathcal{L}(V).
\]
\end{problem}

\begin{proof}
  Denote \( P=\left\{ p(T):p\in \mathcal{P}(\F) \right\}\neq \mathcal{L}(V) \). It's easy to verify that \( P \) is a subspace of \( \mathcal{L}(V) \). \par
  Denote \( d:=\deg\min_T(x)\le n \). We will show that \( \dim P= d < n^2=\dim\mathcal{L}(V) \), which implies \( P \neq \mathcal{L}(V) \). \par
  First, \( I,T,T^2,\dots,T^{d-1} \) is linearly independent in \( \mathcal{L}(V) \), otherwise there exists not-all-zero \( a_0, a_1, \dots, \) \(a_{d-1} \in \F \) such that
  \[ \sum_{k=0}^{d-1}a_kT^k=0\implies\min_T(x)\left|\,\sum_{k=0}^{d-1}a_kx^k\right.\implies a_k=0 \text{ for all }0\le k\le d-1, \]
  which contradicts the choice of \( a_0,a_1,\dots,a_{d-1} \). Hence, \( \dim P\ge d \).\par
  Next, we will show that \( I,T,T^2,\dots,T^{d-1} \) spans \( P \). For any \( p(T)\in P \), we can write \( p(x)=q(x)\min_T(x)+r(x) \) for some \( q,r\in \mathcal{P}(\F) \) such that \( r=0 \) or \( \deg r < d \). Then
  \[ p(T)=q(T)\min{}_T(T)+r(T)=r(T)\in\sspan\left\{ I,T,T^2,\dots,T^{d-1} \right\}. \]
  Hence, \( \dim P\le d \). \par
  In conclusion, \( \dim P=d < n^2=\dim\mathcal{L}(V) \), which implies \( P \neq \mathcal{L}(V) \).
\end{proof}





\begin{problem}[5B-7]
\begin{enumerate}[label=(\alph*)]
  \item Give an example of $S, T \in \mathcal{L}(\F^2)$ such that the minimal polynomial of $ST$ does not equal the minimal polynomial of $TS$.
  \item Suppose $V$ is finite-dimensional and $S, T \in \mathcal{L}(V)$. Prove that if at least one of $S, T$ is invertible, then the minimal polynomial of $ST$ equals the minimal polynomial of $TS$.
\end{enumerate}
\end{problem}

\begin{problem}[5B-16]
Suppose $a_0, \ldots, a_{n-1} \in \F$. Let $T$ be the operator on $\F^n$ whose matrix (with respect to the standard basis) is
\[
\begin{pmatrix}
0 & 0 & 0 & \cdots & 0 & -a_0 \\
1 & 0 & 0 & \cdots & 0 & -a_1 \\
0 & 1 & 0 & \cdots & 0 & -a_2 \\
\vdots & \vdots & \vdots & \ddots & \vdots & \vdots \\
0 & 0 & 0 & \cdots & 0 & -a_{n-2} \\
0 & 0 & 0 & \cdots & 1 & -a_{n-1}
\end{pmatrix}.
\]
Here all entries of the matrix are $0$ except for all $1$'s on the line under the diagonal and the entries in the last column (some of which might also be $0$). Show that the minimal polynomial of $T$ is the polynomial
\[
a_0 + a_1 z + \cdots + a_{n-1}z^{n-1} + z^n.
\]
\end{problem}

\begin{problem}[5B-20]
Suppose $T \in \mathcal{L}(\F^4)$ is such that the eigenvalues of $T$ are $3, 5, 8$. Prove that
\[
(T-3I)^2(T-5I)^2(T-8I)^2=0.
\]
\end{problem}

\begin{problem}[5B-21]
Suppose $V$ is finite-dimensional and $T \in \mathcal{L}(V)$. Prove that the minimal polynomial of $T$ has degree at most $1 + \dim\operatorname{range} T$.
\end{problem}

\begin{problem}[5B-22]
Suppose $V$ is finite-dimensional and $T \in \mathcal{L}(V)$. Prove that $T$ is invertible if and only if
\[
I \in \operatorname{span}(T, T^2, \ldots, T^{\dim V}).
\]
\end{problem}

\begin{problem}[5B-25]
Suppose $V$ is finite-dimensional, $T \in \mathcal{L}(V)$, and $U$ is a subspace of $V$ that is invariant under $T$.
\begin{enumerate}[label=(\alph*)]
  \item Prove that the minimal polynomial of $T$ is a polynomial multiple of the minimal polynomial of the quotient operator $T/U$.
  \item Prove that
  \[
  (\text{minimal polynomial of } T|_U)\times(\text{minimal polynomial of } T/U)
  \]
  is a polynomial multiple of the minimal polynomial of $T$.
\end{enumerate}
\end{problem}

\begin{problem}[5B-29]
Show that every operator on a finite-dimensional vector space of dimension at least two has an invariant subspace of dimension two.
\end{problem}

\end{document}